Existe uma expectativa da parte das famílias amazônicas em relação a natura que se baseia na sustentabilidade econômica das famílias, também acreditamos que existe uma expectativa dos seus consumidores de coerencia da relação da marca com o discurso de sustentabilidade e valorização dos recursos naturais da Amazônia.
A substituição afetaria diretamente as comunidades de forma negativa, e em relação ao seu publico consumidor a substituição é antagônica ao discurso e identidade da marca.
Se ela opar por substituição estaria traindo a confiança dos seus fornecedores e a sua declaração de que o fortalecimento das comunidades de fornecedores esta relacionado a sua estratégia de abastecimento.
Não faz sentido econômico e de imagem a substituição. A única decisão possível é negar a oferta. Dos 12,9 Bilhões em compra 0,19\% do custo de fornecimento se refere ao fornecimento das Comunidade Amazônicas
Não substituir.
